\chapter{Statistics}\label{ch:iv1:statistics}

\begin{ivtrains}
Saying how much a sample tells you: tests, intervals and power, estimators and their errors, multiple testing, the bootstrap and when it breaks, normality, and what a measured Sharpe ratio is worth. Most answers come down to one standard error, computed correctly.
\end{ivtrains}

\begin{question}[$\star$ \iqroles{QR, QT, DS, Analyst, Trader} \iqfirm{any} \iqlevel{junior}]\label{iq:iv1:statistics:1}
A coin gives 60 heads in 100 flips. Test whether it is fair, at the 5\% level, using a normal approximation. Is the result the same with an exact binomial calculation?
\end{question}

\begin{question}[$\star$ \iqroles{QR, QT, DS, Analyst} \iqfirm{any} \iqlevel{junior}]\label{iq:iv1:statistics:2}
Over 400 days a strategy has a mean daily return of $0.1\%$ and a daily standard deviation of $1\%$. Give a 95\% confidence interval for the true mean. What sample size would make the interval half as wide?
\end{question}

\begin{question}[$\star$ \iqroles{QR, QT, DS, Analyst} \iqfirm{any} \iqlevel{junior}]\label{iq:iv1:statistics:3}
A researcher tests ideas of which 10\% are truly good. The test has a 5\% false positive rate and 80\% power. An idea has just passed. What is the chance that it is a dud? Does a p-value of $0.05$ mean a 5\% chance that the null hypothesis is true?
\end{question}

\begin{question}[$\star$ \iqroles{QR, QT, DS, Analyst} \iqfirm{any} \iqlevel{junior}]\label{iq:iv1:statistics:4}
Why does the sample variance divide by $n - 1$ instead of $n$? Show that it makes the estimator unbiased. What does the plain estimator with $n$ average for samples of size 5?
\end{question}

\begin{question}[$\star$ \iqroles{QR, QT, DS, Analyst} \iqfirm{any} \iqlevel{junior}]\label{iq:iv1:statistics:5}
You test 20 independent trading signals, none of which works, each at the 5\% level. How likely is it that one or more pass? What level per test would keep that chance at 5\%?
\end{question}

\begin{question}[$\star$ \iqroles{QR, QT, Analyst, Trader} \iqfirm{hedge fund} \iqlevel{junior}]\label{iq:iv1:statistics:6}
A thousand managers have no skill, and each beats the market in any one year with a 50\% chance, independently. After five years, how many on average have beaten it every single year? What does an investor who screens for a clean five-year record learn?
\end{question}

\begin{question}[$\star$ \iqroles{QR, QT, DS, Analyst} \iqfirm{any} \iqlevel{junior}]\label{iq:iv1:statistics:7}
A strategy is up on 15 of 20 days. With no assumption about the size or the shape of returns, test whether its median daily return is zero. Give the exact two-sided p-value.
\end{question}

\begin{question}[$\star$ \iqroles{QR, QT, DS} \iqfirm{any} \iqlevel{junior}]\label{iq:iv1:statistics:8}
Algorithm A produced the values 5, 7 and 9, and algorithm B produced 1, 2 and 3. With no assumption about distributions, test whether A gives higher values, using a permutation test. What is the exact one-sided p-value? Why can no sample of this size reach below it?
\end{question}

\begin{question}[$\star\star$ \iqroles{QR, QT, DS, Analyst} \iqfirm{any} \iqlevel{junior}]\label{iq:iv1:statistics:9}
A sample of 5 has mean 1 and standard deviation 1. Give 95\% intervals for the true mean, one using the normal value $1.96$ and one using the correct $t$ value, given that the 97.5\% point of a $t$ distribution with 4 degrees of freedom is $2.776$. Why do they differ, and when does the difference stop mattering?
\end{question}

\begin{question}[$\star\star$ \iqroles{QR, QT, DS, Analyst} \iqfirm{any} \iqlevel{junior}]\label{iq:iv1:statistics:10}
Two series have a sample correlation of $0.5$ over 103 points. Give the rule of thumb for the standard error of a correlation near zero, then a proper 95\% interval for this one using Fisher's transformation. Why is the interval not symmetric around $0.5$?
\end{question}

\begin{question}[$\star\star$ \iqroles{QR, QT, DS, Risk} \iqfirm{any} \iqlevel{junior}]\label{iq:iv1:statistics:11}
In 100 trading days a system has had no failures. What upper bound can you give, with 95\% confidence, on its failure chance per day? Give the exact number and the rule of thumb. What is the naive estimate, and why is it not useful?
\end{question}

\begin{question}[$\star\star$ \iqroles{QR, QT, Analyst, Trader} \iqfirm{hedge fund} \iqlevel{junior}]\label{iq:iv1:statistics:12}
A strategy has a true annual Sharpe ratio of $0.5$. How many years of daily returns do you need before its measured Sharpe ratio is, on average, two standard errors above zero? What about true Sharpe ratios of 1 and 2? What does this say about the length of track records?
\end{question}

\begin{question}[$\star\star$ \iqroles{QR, QT, DS, Analyst} \iqfirm{hedge fund, market maker} \iqlevel{junior}]\label{iq:iv1:statistics:13}
A signal adds $0.05\%$ a day to a return with a daily standard deviation of $1\%$. For a two-sided test at the 5\% level with 80\% power, how many days of data do you need? Give the formula and the result in years. How does the answer change if the effect is twice as large?
\end{question}

\begin{question}[$\star\star$ \iqroles{QR, QT, DS} \iqfirm{any} \iqlevel{junior}]\label{iq:iv1:statistics:14}
Ten tests produce the p-values $0.001$, $0.008$, $0.012$, $0.041$, $0.042$, $0.06$, $0.074$, $0.205$, $0.212$ and $0.216$. How many are rejected under Bonferroni at a family level of $0.05$? How many under Benjamini--Hochberg, with a false discovery rate of $0.05$? What does each one control?
\end{question}

\begin{question}[$\star\star$ \iqroles{QR, QT, DS} \iqfirm{any} \iqlevel{junior}]\label{iq:iv1:statistics:15}
Waiting times follow a Gamma distribution with shape $k$ and scale $\theta$, whose mean is $k\theta$ and variance is $k\theta^2$. A sample has mean 4 and variance 8. Give the method-of-moments estimates. Why would you still prefer the maximum-likelihood estimates, if you can compute them?
\end{question}

\begin{question}[$\star\star$ \iqroles{QR, QT, DS} \iqfirm{any} \iqlevel{junior}]\label{iq:iv1:statistics:16}
You have $n$ independent exponential waiting times with rate $\lambda$. Derive the maximum-likelihood estimate of the rate. Is it unbiased? Give the unbiased version, and its value when $n = 5$ and the sample mean is 2.
\end{question}

\begin{question}[$\star\star$ \iqroles{QR, QT, DS} \iqfirm{any} \iqlevel{junior}]\label{iq:iv1:statistics:17}
Daily counts are Poisson with mean $\lambda$, and you have $n$ days. Show that the sample mean reaches the lower bound on the variance of an unbiased estimator. To estimate the chance of a day with no events, compare the fraction of zero days with $e^{-\bar{X}}$, at $\lambda = 2$, in asymptotic variance. Which is better, and what does the better one assume?
\end{question}

\begin{question}[$\star\star$ \iqroles{QR, QT, DS, QD} \iqfirm{any} \iqlevel{junior}]\label{iq:iv1:statistics:18}
You draw a bootstrap sample of size $n$ with replacement from $n$ data points. What fraction of the original points does it contain on average, for large $n$? What does the bootstrap estimate, and in plain words why does it work?
\end{question}

\begin{question}[$\star\star$ \iqroles{QR, QT, DS, Risk} \iqfirm{any} \iqlevel{junior}]\label{iq:iv1:statistics:19}
Daily returns follow an AR(1) with coefficient $\rho = 0.5$. By what factor does the variance of the sample mean exceed the value computed by assuming independence? What does a plain bootstrap report, and what should you do instead?
\end{question}

\begin{question}[$\star\star$ \iqroles{QR, QT, DS, Risk} \iqfirm{any} \iqlevel{junior}]\label{iq:iv1:statistics:20}
A thousand daily returns have a skewness of $-0.5$ and an excess kurtosis of 3. Give the standard errors of those two numbers under normality. Compute the Jarque--Bera statistic, and say what it is compared with. What is the weakness of this test for financial data?
\end{question}

\begin{question}[$\star\star$ \iqroles{QR, QT, Risk, Trader, Analyst} \iqfirm{any} \iqlevel{junior}]\label{iq:iv1:statistics:21}
A daily return has a standard deviation of $1\%$, and today it is $-5\%$. If returns are normal, how often should you see such a day, in years? How often if returns follow a Student $t$ with 3 degrees of freedom, rescaled to the same standard deviation? What do you conclude?
\end{question}

\begin{question}[$\star\star$ \iqroles{QR, QT, DS} \iqfirm{any} \iqlevel{junior}]\label{iq:iv1:statistics:22}
A die is rolled 120 times, and the faces 1 to 6 come up 30, 18, 20, 22, 16 and 14 times. Run a chi-square test of fairness, given that the 95\% point with 5 degrees of freedom is $11.07$. A casual look suggests that the face 1 is favoured. Would you test that one face alone, and how?
\end{question}

\begin{question}[$\star\star$ \iqroles{QR, QT, DS, Trader} \iqfirm{any} \iqlevel{junior}]\label{iq:iv1:statistics:23}
A string of 100 up and down days contains 30 runs, where a run is a maximal block of the same sign. If the days are independent and fair, what mean and standard deviation does the number of runs have? What does 30 suggest? What would 65 suggest?
\end{question}

\begin{question}[$\star\star$ \iqroles{QR, QT, DS, LLE} \iqfirm{market maker, hedge fund} \iqlevel{junior}]\label{iq:iv1:statistics:24}
Two execution algorithms each trade the same 100 orders. Their slippages have a standard deviation of 10 basis points each, a correlation of $0.9$ across orders, and a mean difference of 1 basis point. What is the $z$-score if you compare them as two independent samples? If you compare them as pairs? Which one is right, and what does this suggest for designing a test?
\end{question}

\begin{question}[$\star\star$ \iqroles{QR, QT, DS, Analyst} \iqfirm{any} \iqlevel{junior}]\label{iq:iv1:statistics:25}
Two independent estimates are 0 and 3, each with a standard error of 1. Their 95\% intervals overlap. Is the difference significant at the 5\% level? Find the general rule relating overlap to significance.
\end{question}

\begin{question}[$\star\star$ \iqroles{QR, QT, DS, Risk} \iqfirm{any} \iqlevel{junior}]\label{iq:iv1:statistics:26}
For ten points, $x$ runs from 1 to 10, and $y = x$, except that the last $y$ is 100. Compute the Pearson and the Spearman correlations. Which gives the fairer summary, and when would you worry about using each of them?
\end{question}

\begin{question}[$\star\star\star$ \iqroles{QR, QT, DS} \iqfirm{hedge fund} \iqlevel{junior, experienced}]\label{iq:iv1:statistics:27}
You estimate the variance of a normal distribution from $n$ independent draws, using $\frac{1}{c}\sum (X_i - \bar{X})^2$. Find the divisor $c$ with the smallest mean squared error. For $n = 10$, give the mean squared error, in units of $\sigma^4$, for $c = n - 1$, $n$ and the best $c$. How does that square with unbiasedness being a virtue?
\end{question}

\begin{question}[$\star\star\star$ \iqroles{QR, QT, DS} \iqfirm{hedge fund} \iqlevel{junior, experienced}]\label{iq:iv1:statistics:28}
You must choose between two exact hypotheses from $n$ independent draws: the distribution is $N(0, 1)$, or it is $N(1, 1)$. What is the best test, and why? How large must $n$ be for both error rates to be at most 5\%, if they are set equal?
\end{question}

\begin{question}[$\star\star\star$ \iqroles{QR, QT, DS, Analyst} \iqfirm{hedge fund, market maker} \iqlevel{junior, experienced}]\label{iq:iv1:statistics:29}
You run an experiment on a null effect. After each of 10 equally spaced batches, you compute a $z$-score on all data so far, and stop as soon as $|z| > 1.96$. About what is your real false positive rate? What happens with no limit on looks? What cut-off per look would restore 5\%, and what is the method called?
\end{question}

\begin{question}[$\star\star\star$ \iqroles{QR, QT, DS} \iqfirm{any} \iqlevel{junior, experienced}]\label{iq:iv1:statistics:30}
You draw $n$ independent values from a uniform on $[0, \theta]$, and you want a confidence interval for $\theta$ from the sample maximum. Show that a plain bootstrap of the maximum fails, by finding the probability that the bootstrap maximum equals the sample maximum, for large $n$. What does the true sampling distribution look like, and what would you use?
\end{question}

\begin{question}[$\star\star\star$ \iqroles{QR, QT, Analyst, Risk} \iqfirm{hedge fund} \iqlevel{junior, experienced}]\label{iq:iv1:statistics:31}
A strategy has a monthly Sharpe ratio of $s = 0.3$ measured over $n = 120$ months. Derive the standard error of the estimate under normality. How much wider does the interval get if monthly returns have a skewness of $-1$ and a kurtosis of 6, given that the variance becomes $\bigl(1 - \gamma_3 s + \frac{\gamma_4 - 1}{4}s^2\bigr)/n$? Why does it get wider?
\end{question}

\begin{question}[$\star\star\star$ \iqroles{QR, QT, Analyst, Risk} \iqfirm{hedge fund} \iqlevel{junior, experienced}]\label{iq:iv1:statistics:32}
Monthly returns have a Sharpe ratio of $0.3$ and first-order autocorrelation of $0.2$, following an AR(1). The usual rule multiplies by $\sqrt{12}$. What is the correct annual Sharpe ratio? What happens with negative autocorrelation? What kinds of strategies display positive autocorrelation in their reported returns, and why?
\end{question}

\begin{question}[$\star\star\star$ \iqroles{QR, QT, DS} \iqfirm{hedge fund} \iqlevel{junior, experienced}]\label{iq:iv1:statistics:33}
Two strategies have the same true annual Sharpe ratio and a correlation of $0.8$, and you have 5 years of daily returns on each. What is the standard error of the difference of the two measured Sharpe ratios? How does that compare with treating them as independent? What does this imply about choosing between similar strategies?
\end{question}

\begin{question}[$\star\star\star$ \iqroles{QR, QT, Risk, DS} \iqfirm{bank, hedge fund} \iqlevel{junior, experienced}]\label{iq:iv1:statistics:34}
You estimate the 1\% quantile of a daily return, value at risk, from 1000 days of normal returns. Derive the standard error of the sample quantile, and compare it with the error of $2.326\hat{\sigma}$, in units of $\sigma$. Why is the first larger? When should you still choose the first?
\end{question}

\begin{question}[$\star\star\star$ \iqroles{QR, QT, Risk, DS} \iqfirm{any} \iqlevel{junior, experienced}]\label{iq:iv1:statistics:35}
A sample of 101 returns has a standard deviation of $2\%$. What is the standard error of that estimate if returns are normal? If they have a kurtosis of 6? Derive the rule. What does this do to the usual chi-square interval?
\end{question}

\begin{question}[$\star\star\star$ \iqroles{QR, QT, DS} \iqfirm{hedge fund} \iqlevel{junior, experienced}]\label{iq:iv1:statistics:36}
You measure a correlation of $0.02$ between a signal and next-day returns. Half of the signal's variance is measurement noise, and three-quarters of the variance of the day's return is noise unrelated to anything you could forecast. What is the correlation between the true signal and the forecastable part of returns? Derive the formula. What does this say about raw correlations in finance?
\end{question}

\begin{question}[$\star\star\star$ \iqroles{QR, QT, DS, Risk} \iqfirm{hedge fund, market maker} \iqlevel{junior, experienced}]\label{iq:iv1:statistics:37}
An asset's return has an annual volatility of 20\%, and you have 10 years of daily data. Find the standard errors of the annual mean and of the annual volatility, in percentage points. How does each change if you sample monthly, or by the minute? What practical conclusion follows, and what prevents you from pushing the volatility result without limit?
\end{question}

\begin{question}[$\star\star\star\star$ \iqroles{QR, QT, DS} \iqfirm{hedge fund} \iqlevel{experienced}]\label{iq:iv1:statistics:38}
You observe $X \sim N(\theta, I_d)$ in $d$ dimensions and estimate the mean $\theta$ with total squared error as the loss. Show that the James--Stein estimator $\bigl(1 - (d - 2)/|X|^2\bigr)X$ beats $X$ itself for every $\theta$ when $d \ge 3$. Give its risk at $\theta = 0$ for $d = 10$. Why is it strange, and how would you use it on a set of asset mean returns?
\end{question}

\begin{question}[$\star\star\star\star$ \iqroles{QR, QT, DS} \iqfirm{hedge fund} \iqlevel{experienced}]\label{iq:iv1:statistics:39}
A sample of $n$ comes from a Cauchy distribution with centre $\theta$, with density $\frac{1}{\pi(1 + (x - \theta)^2)}$. Compute the Fisher information per observation. Compare the asymptotic variances of the maximum-likelihood estimate, the sample median and the sample mean. How would you compute the first in practice, given that its likelihood can have several local peaks?
\end{question}

\begin{question}[$\star\star\star\star$ \iqroles{QR, QT, DS} \iqfirm{hedge fund} \iqlevel{experienced}]\label{iq:iv1:statistics:40}
Prove that the Benjamini--Hochberg procedure at level $\alpha$ controls the false discovery rate at $\alpha m_0/m$ when the $m$ p-values are independent, with $m_0$ true nulls, whose p-values are uniform. Hint: for one true null, ask what the number of rejections would be if its p-value were set to 0.
\end{question}
